\documentclass[11pt]{article}
\usepackage[margin=1in]{geometry}
\usepackage{amsmath,amssymb,amsthm}
\usepackage{enumitem}
\usepackage{algorithm}
\usepackage{algpseudocode}

\newtheorem{theorem}{Theorem}
\newtheorem{proposition}[theorem]{Proposition}
\newtheorem{corollary}[theorem]{Corollary}
\newtheorem{lemma}[theorem]{Lemma}
\theoremstyle{definition}
\newtheorem{definition}[theorem]{Definition}
\theoremstyle{remark}
\newtheorem{remark}[theorem]{Remark}

\newcommand{\eps}{\tilde{\epsilon}}
\newcommand{\phist}{\tilde{\varphi}^*}
\newcommand{\Z}{\mathbb{Z}}

\title{Finishing-Feasibility at the Last Two Levels\\
for Spoof Lehmer Enumeration}
\author{G.~Molnar and G.~Singh}
\date{\today~(Draft)}

\begin{document}
\maketitle

\begin{abstract}
We sharpen the bounds-propagation algorithm for enumerating spoof Lehmer
factorizations by replacing the last two levels of recursive search with
direct algebraic solves.  At length $s = r - 1$, the residual factor is
given by a closed-form rational; at length $s = r - 2$, the residual pair
is determined by divisor pairs of a single explicit integer.  The
sharpening is specific to the defining equation and is especially
beneficial at small $k$, where the bounds-propagation geometry is loose
and the recursion is correspondingly deep.  We also note that the
$s = r - 2$ identity specializes to the Hasanalizade extension equation
when the prefix is a plus-seed, unifying the small-$k$ completeness
enumeration with the descent framework of~\cite{MFT}.
\end{abstract}

\section{Introduction}

In~\cite{MS26} we introduced spoof Lehmer factorizations---multisets
$F = \{x_1, \ldots, x_r\}$ of positive integers satisfying
$k \cdot \phist(F) = \eps(F) - 1$ for some $k \in \Z_{>0}$, where
$\eps(F) = \prod x_i$ and $\phist(F) = \prod(x_i - 1)$---and
enumerated all such factorizations with $r \le 6$.  Our enumeration
proceeded by \emph{bounds propagation}: building up sorted partial
factorizations $F_0 = (x_1, \ldots, x_s)$ with $s < r$, and pruning
candidates for $x_{s+1}$ using the monotone ratio
\begin{equation}\label{eq:ratio}
  k(F) = \frac{\eps(F) - 1}{\phist(F)},
\end{equation}
which is decreasing in each argument.  Specifically, the lower and upper
bounds
\[
  L(F_0) = \frac{\eps(F_0) - [s = r]}{\phist(F_0)},
  \qquad
  U(F_0) = \frac{\eps(F_0') - 1}{\phist(F_0')},
\]
where $F_0'$ fills the remaining $r - s$ slots with the smallest legal
next factor, sandwich the achievable values of $k$ for any completion
of~$F_0$.  The recursion terminates whenever $U < k$ or $L > k$.

This algorithm successfully enumerates all $r \le 6$ spoof Lehmer
factorizations.  For $r = 7$, the algorithm is in principle correct but
empirically stalls at $k = 2$.  The reason is structural: at small~$k$,
the shared limit $\lim_{a \to \infty} L(F_0, a) = \lim_{a \to \infty}
U(F_0, a) = \eps(F_0)/\phist(F_0)$ approaches the target from the
\emph{constraint} side of $k$, so the termination inequality $U < k$ only
fires after many candidate factors have been explored.  At large~$k$,
partial states are pruned aggressively; at $k = 2$, the surviving
frontier balloons.

In this note we sharpen the algorithm by replacing the last two levels
of recursion with direct solves.

\section{Finishing-feasibility at $s = r - 1$}

Let $F_0 = (x_1, \ldots, x_{r-1})$ be a partial factorization with
$\eps(F_0) = E$ and $\phist(F_0) = P$.  A single factor~$x_r$ completes
$F_0$ to a $k$-Lehmer factorization iff
$k P (x_r - 1) = E x_r - 1$.  Solving:
\begin{equation}\label{eq:r-1-finish}
  x_r \;=\; \frac{kP - 1}{kP - E}.
\end{equation}

\begin{proposition}[$s = r - 1$ finishing-feasibility]\label{prop:rm1}
A partial factorization $F_0$ of length $r - 1$ extends to a $k$-Lehmer
factorization if and only if $kP - E > 0$, $kP - E$ divides $kP - 1$,
and the quotient $x_r = (kP - 1)/(kP - E)$ satisfies $x_r \ge x_{r-1}$.
The extension, when it exists, is unique.
\end{proposition}

The entire check requires one integer division and one comparison.  No
candidate iteration is needed.

\section{Finishing-feasibility at $s = r - 2$}

Let $F_0 = (x_1, \ldots, x_{r-2})$ with $\eps(F_0) = E$ and
$\phist(F_0) = P$.  We seek a pair $(a, b)$ with $a \le b$ and
$a \ge x_{r-2}$ completing $F_0$ to a $k$-Lehmer factorization.  The
defining equation expands to
\[
  kP(a - 1)(b - 1) \;=\; Eab - 1,
\]
which rearranges to
\[
  (kP - E)\,ab \;-\; kP(a + b) \;+\; (kP + 1) \;=\; 0.
\]
Write $A = kP - E$ and $B = kP$.  Multiplying by $A$ and completing the
product on $a$ and $b$ yields
\begin{equation}\label{eq:finish}
  (A a - B)(A b - B) \;=\; B E - A.
\end{equation}

\begin{proposition}[$s = r - 2$ finishing-feasibility]\label{prop:rm2}
Let $F_0$ be a partial factorization of length $r - 2$, with $A = kP - E$
and $B = kP$.  A pair $(a, b)$ with $a \le b$ and $a \ge x_{r-2}$ extends
$F_0$ to a $k$-Lehmer factorization if and only if $A > 0$ and there is
a divisor pair $(d_1, d_2)$ of $BE - A$ with $d_1 \le d_2$ such that
\[
  a = \frac{d_1 + B}{A},
  \qquad
  b = \frac{d_2 + B}{A},
\]
are integers satisfying $x_{r-2} \le a \le b$.
\end{proposition}

The check requires one factorization of the integer $BE - A$, followed
by one trial of each divisor pair.  The cost is dominated by the
factorization step.

\begin{remark}[Relationship to the descent formula]
When $F_0$ is a \emph{plus-seed}---i.e., $k P = E + 1$, so $A = 1$ and
$B = E + 1$---Equation~\eqref{eq:finish} becomes
$(a - E - 1)(b - E - 1) = (E + 1) E - 1 = E^2 + E - 1$, which is exactly
Hasanalizade's extension equation~\cite{H26} specialized to the Lehmer
side.  More generally, the $s = r - 2$ finishing-feasibility identity
subsumes the Hasanalizade descent equation: the descent formula handles
the case $A = 1$; Proposition~\ref{prop:rm2} handles arbitrary $A \ge 1$.
\end{remark}

\section{Algorithm}

The sharpened bounds-propagation algorithm replaces the last two levels
of recursion with direct solves.

\begin{algorithm}[H]
\caption{Sharpened bounds propagation for $k$-Lehmer enumeration}
\label{alg:sharp}
\begin{algorithmic}[1]
\Procedure{Enumerate}{$F_0$, $k$, $r$}
\State $s \gets |F_0|$
\If{$s = r - 1$}
    \State solve~\eqref{eq:r-1-finish}; emit if integer and valid
\ElsIf{$s = r - 2$}
    \State factor $BE - A$; emit every divisor pair giving valid $(a, b)$
\Else
    \State apply bounds propagation as in~\cite{MS26}
    \For{each candidate $x_{s+1}$ not pruned by $L, U$ or Lehmer's congruence}
        \State \Call{Enumerate}{$F_0 \cup \{x_{s+1}\}$, $k$, $r$}
    \EndFor
\EndIf
\EndProcedure
\end{algorithmic}
\end{algorithm}

The last two levels change from "explore a potentially large frontier of
candidate factors" to "solve one equation and verify."  Both finishing
checks are dominated by the $s = r - 2$ factorization; everything above
$s = r - 2$ is unchanged from the original algorithm.

\begin{remark}[Expected impact at $k = 2$]
At $k = 2$, the bounds-propagation recursion's next-factor loop runs
longest precisely at the deepest levels $s = r - 1$ and $s = r - 2$,
because the gap $\min(k - L, U - k)$ is tightest there but the bound
$U$ decays slowest as $a$ grows.  Replacing these levels with direct
solves eliminates the loop at exactly the depths where the original
algorithm spent the most time.
\end{remark}

\section{Implementation notes}

Algorithm~\ref{alg:sharp} is straightforward to implement on top of the
existing bounds-propagation enumerator.  The only external dependency is
an integer factorization routine for $BE - A$; for the $r = 7$ case of
interest, this integer has at most around~70 digits and is within reach
of standard methods (sympy's Pollard rho plus ECM, or yafu for larger
inputs).  The divisor enumeration runs in time
$O(\tau(BE - A) \log(BE - A))$, which is negligible compared to the
factoring step.

Lehmer's congruence pruning (no factor of a $k$-Lehmer factorization is
$\equiv 1 \pmod{p}$ for any other factor~$p$) continues to apply at
levels $s < r - 2$ and remains useful there.  At $s = r - 1$ and
$s = r - 2$, the congruence is subsumed by the direct-solve conditions:
factors violating the congruence necessarily fail the divisibility or
integrality checks.

\begin{thebibliography}{9}

\bibitem{H26}
E.~Hasanalizade, \emph{Some remarks on a paper of Molnar and Singh},
Integers \textbf{26} (2026).

\bibitem{MFT}
G.~Molnar, \emph{The extension graph of spoof Lehmer factorizations is a
forest}, in preparation (2026).

\bibitem{MS26}
G.~Molnar and G.~Singh, \emph{Spoof Lehmer factorizations}, Integers
\textbf{26} (2026), \#A33.

\end{thebibliography}

\end{document}
